\begin{tomtat}
\begin{itemize}
  \item Cái sổ mà bảy chương Phần~IV lập gồm chín bài báo, không phải sáu, và nó
không có giao. Sáu bài nêu cùng một thứ đang thiếu, là một dụng cụ đo độ trung
thực đã được kiểm định, và nêu từ sáu hướng; ba bài không nêu, và ba chỗ trống
của chúng cũng khác nhau, một bài muốn một đại lượng tính được, một bài muốn mô
hình khái quát hóa tốt hơn, một bài muốn khám phá nhân quả tốt hơn. Lấy giao của
chín tập hợp mà ba trong số đó không chứa phần tử chung nào thì được tập rỗng.
  \item Hỏi cái sổ một câu hẹp hơn thì được một con số khác, và câu ấy là: phần
hạn chế hoặc phần công việc tương lai \emph{của chính bài ấy} có đòi cái dụng cụ
đang thiếu hay không. Câu ấy cho hai trên chín, và chỉ một bài gọi tên được dụng
cụ, là các đường cong deletion và insertion trên từng cặp token. Bảy dòng còn lại
phần lớn đòi một bản tốt hơn của đúng thứ bài ấy vốn đã dựng; dòng gần nhất với
việc đòi một dụng cụ đo là dòng đòi một chỉ số đặt nền trên cảm nhận của người,
mà tiêu chí ấy nằm ở dòng tính hợp lý của bảng~\ref{tab:ch01-bon-quan-he} chứ
không ở dòng độ trung thực. \textbf{Khoảng cách giữa sáu và hai là khoảng cách
giữa chẩn đoán và đơn thuốc}, và nó là kết quả chính của chương.
  \item Ba tài liệu ngoài danh mục 32 bài lẽ ra đã khép được khoảng trống và
không bài nào khép. Bài siêu đánh giá lấy chính các chỉ số làm đối tượng và
tuyên bố ở trang 5 rằng không thể xác định tính hợp lệ của một bộ ước lượng vì
việc ấy đòi nhãn ground truth, rồi chuyển sang phát biểu các chế độ hỏng, tức
đúng nước đi mà mục~\ref{sec:ch08-dieu-kien-can} mô tả ở một tầng dưới. Bài tổng
quan mà ba trong bốn bài lõi của Phần~IV cùng dẫn là một bản kiểm kê không chạy
thí nghiệm nào, và câu hướng tới tương lai của nó đòi kiểm định các
\emph{phương pháp} bằng chính bộ công cụ nó vừa kiểm kê. Còn bài chuẩn thì đã in
sẵn số liệu để hỏi câu ấy và đọc nó thành chuyện khác. Trong bảng của nó, ở hàng \code{monotonicity}, cả bảy cột nằm gọn giữa
0.517 và 0.562 còn cột mốc ngẫu nhiên đứng ở 0.525, tức không đứng cuối; ở hàng
\code{GT-Shapley}, cũng bảy cột ấy trải từ $-0.127$ tới 0.930 và cột mốc đứng ở
0.004. Phương pháp bị chấm tệ nhất ở hàng sau là phương pháp được chấm tốt nhất ở
hàng trước, và chính bài giải thích vì sao: phương pháp ấy chọn đặc trưng theo
đúng thao tác mà \code{monotonicity} kiểm. Bảng ấy không khép được khoảng trống
vì đáp án mà \code{GT-Shapley} đối chiếu là giá trị Shapley chính xác, tức một
định nghĩa attribution đã được chọn rồi tính không sai số, nên nó so một lời giải
thích với một lời giải thích khác chứ không so lời giải thích với mô hình.
  \item Phép đối chiếu đã chạy được ở đúng hai chỗ, và cái chung của hai chỗ ấy
không phải sự sắc sảo mà là cái giá bằng không: một bên vế phải lộ ra vì thiết kế
nhiệm vụ ép nó lộ, một bên vì khái niệm đã được gán nhãn sẵn cùng dữ liệu cho một
mục đích khác. Ở phía attribution đặc trưng thì vế phải phải dựng lấy, tức giá
khác không, và ở đó phép đối chiếu chưa được chạy. Nguyên tắc dựng của bài neo
không gắn với chuỗi suy luận, nên phía kia dùng được nó.
  \item Năm chương liền, mỗi chương đọc một thứ khác nhau, rút ra năm điều kiện
đặt lên một dụng cụ đo tương lai, và bảng~\ref{tab:ch18-nam-dieu-kien} xếp cả năm:
mô tả chỗ lệch thay vì xác nhận sự trùng khớp, công bố tập can thiệp, báo phép ép
cạnh phép đọc tên, khai ra dụng cụ nằm dưới đã được kiểm định bằng gì, và nói rõ
nó hỏi câu nào trong hai câu. Không điều kiện nào được phát biểu bởi một bài báo
nào; cả năm là phần còn lại sau khi một chương đọc xong bài của nó.
  \item Sách không nói các chỉ số sai. Chưa kiểm định không phải là sai: một dụng
cụ chưa được đối chiếu vẫn có thể đang đo đúng thứ ghi trên mặt đồng hồ, và cái
không có là lý do để tin như vậy. Sách cũng không nói lời giải thích vô ích, và
không quy khoảng trống cho lỗi của ai: mỗi bài đi một bước hợp lý, và khoảng
trống là thứ chỉ hiện ra khi xếp mười bảy chương cạnh nhau.
\end{itemize}
\end{tomtat}

\cauhoi
\begin{cauhoids}
  \item Mục~\ref{sec:ch18-so-va-chu-giao} lập luận rằng lấy giao của chín tập hợp
mà ba trong số đó không giao với phần còn lại thì được tập rỗng. Hãy nêu một cách
đọc chữ \enquote{giao} khác, dưới đó lời chỉ dẫn ban đầu vẫn có nghĩa, rồi cho
biết dưới cách đọc ấy thì ba bài không giao đóng vai trò gì.
  \item Bảng~\ref{tab:ch18-so-chan-doan} có hai cột trả lời hai câu hỏi khác
nhau. Hãy chọn một dòng mà hai cột lệch nhau nhiều nhất và viết ra, bằng lời của
bạn, một câu mà bài ấy \emph{đã có thể} viết trong phần công việc tương lai để
hai cột khớp nhau, rồi cho biết vì sao viết câu ấy lại khó với chính tác giả của
bài.
  \item Mục~\ref{sec:ch18-ba-bai-le-ra-da-dong} đọc nước đi của bài siêu đánh giá
là cùng nước đi với mục~\ref{sec:ch08-dieu-kien-can}, chỉ ở một tầng khác. Hãy
kiểm cách đọc ấy: nêu một điều kiện cần mà một \emph{chỉ số} phải thỏa, tương tự
như những \tn{phép thử ngẫu nhiên hóa}{randomization test} đặt lên một bản đồ nổi bật, rồi cho biết một
chỉ số thỏa điều kiện ấy thì ta được phép kết luận gì.
  \item Ở hàng chỉ số đơn điệu của bảng mà mục~\ref{sec:ch18-ba-bai-le-ra-da-dong}
đọc, cột ngẫu nhiên đứng ở 0.525 và một phương pháp đã công bố đứng ở 0.517. Hãy
cho biết dữ kiện ấy một mình chứng minh được điều gì và không chứng minh được
điều gì, rồi nêu một phép đo thêm mà bạn sẽ chạy trước khi kết luận bất cứ điều
gì về chỉ số ấy.
  \item Bài chuẩn giải thích rằng phương pháp thắng ở hàng chỉ số đơn điệu vì nó
chọn đặc trưng theo đúng thao tác mà chỉ số ấy kiểm. Hãy chọn một chỉ số khác
trong bốn họ mà chương~\ref{ch:do-do-trung-thuc} liệt kê và mô tả, bằng hai hoặc
ba câu, một phương pháp attribution được thiết kế riêng để thắng nó; rồi cho biết
phương pháp ấy có nhất thiết là một phương pháp tồi hay không.
  \item Hình~\ref{fig:ch18-hinh-dang-khoang-trong} có ba cột và hai cột đã chạy
được phép đối chiếu. Hãy nêu một loại lời giải thích thứ tư, không nằm trong ba
cột ấy, rồi điền hai ô của nó: vế phải lấy ở đâu, và cái giá của việc lấy nó là
gì. Nếu bạn cho rằng loại ấy rơi vào cột thứ ba thì nói vì sao.
  \item Bảng~\ref{tab:ch18-nam-dieu-kien} nói điều kiện thứ năm không đòi công bố
thêm thứ gì. Hãy chọn một trong tám chỉ số mà mục~\ref{sec:ch09-ket-qua} liệt kê
và một trong các họ chỉ số của chương~\ref{ch:do-do-trung-thuc}, rồi cho biết
dưới điều kiện thứ năm thì mỗi bên phải khai gì, và hai lời khai ấy có giống nhau
không.
  \item Mục~\ref{sec:ch18-phep-do-chua-ai-chay} xếp bốn phép đo theo giá và nói
cái rẻ nhất chưa ai làm. Hãy viết ra, dưới dạng bốn hoặc năm bước, quy trình để
chạy phép đo rẻ nhất ấy trên số liệu đã công bố, rồi cho biết bước nào trong số
đó là bước mà một người đọc không có quyền truy cập gì thêm vẫn làm được.
  \item Đọc mục 2.2 và mục 3.1 trong \enquote{The Meta-Evaluation Problem in
Explainable AI: Identifying Reliable Estimators with MetaQuantus} của Hedström và
cộng sự~\autocite{mmetaquantus}. Ở mục 2.2, bài chia các đại lượng của bài toán
thành hai nhóm, nhóm có sẵn nhãn ground truth và nhóm không, rồi lập luận rằng
chỗ hỏng nằm ở nhóm thứ hai. Hãy cho biết cách dựng ground truth của bài neo, mà
mục~\ref{sec:ch09-dung-ground-truth} mô tả, chuyển được thứ gì từ nhóm thứ hai
sang nhóm thứ nhất, và thứ gì thì không.
  \item Mục~\ref{sec:ch18-khoang-trong-phat-bieu-ra} nói sách không khẳng định
các chỉ số sai, mà chỉ khẳng định chúng chưa được cấp tư cách bằng chứng. Hãy
dựng một tình huống trong đó phân biệt ấy đổi một quyết định thực tế: nêu người
quyết định là ai, họ định dùng điểm số vào việc gì, và họ làm khác đi thế nào
dưới hai cách đọc.
\end{cauhoids}
